% !TeX root = main_fps.tex
\section{Methodology}

The Field Particle System (FPS) is a particle methodology
in which particle interactions are governed by locally
accessible quantities evaluated at the present instant.
Particle interactions are represented through geometric
contact states, and forces arise directly from the
evolving interaction geometry. State updates are
performed through a source-owned formulation in which
particles evaluate interactions independently and modify
only their own state variables. Consequently, force
evaluation, acceleration accumulation, and state
evolution remain locally owned throughout the simulation.

FPS is source-owned, meaning that particle interactions are evaluated from the 
perspective of the source particle and state updates are restricted to quantities 
owned by the evaluating particle. Shared-state modification does not occur 
during pair interactions.

The interaction model is purely geometric, with
forces arising from geometric relationships between
particles. Particle separation, overlap depth, and
contact orientation define the interaction state
from which force and particle motion are
subsequently determined.

The governing dynamics are instantaneous, meaning that
particle interactions are evaluated from quantities
available at the present instant. Previous states and
future states do not participate directly in the
governing dynamics.

The contact model permits particles to occupy a common
geometric region during interaction. The resulting
penetration depth is treated as a physical interaction
state rather than a numerical error requiring immediate
correction.

The methodology does not employ energy as a governing
state variable. Energy may be computed as a diagnostic
quantity but does not directly participate in the
governing dynamics.

The present implementation employs normalized state variables 
rather than dimensional physical quantities. Variables such as 
velocity, force, acceleration, momentum, and particle size are 
represented using relative values that typically lie within the 
interval $[0,1]$. Particle behavior is therefore governed by the 
relative relationships between quantities rather than their absolute 
physical magnitudes. The present work focuses on the geometric 
and dynamical behavior of the methodology independent of a 
particular physical scale.

The methodology of FPS is founded upon a set of
fundamental properties that define the governing
dynamics and interaction model. The following sections 
describe the geometric quantities and 
dynamic relationships used to evolve the system 
through time.


\subsection{FPS Dynamics}

Particles are represented by circular boundaries
defined by a center location $\mathbf{x}_i$ and an
effective radius $r_i$. Particle interactions occur
when the boundaries of two particles interpenetrate,
establishing an interaction state. The resulting penetration 
depth serves as the fundamental interaction quantity 
from which force, acceleration, and particle motion 
are subsequently determined.

Consider two particles having center locations
$\mathbf{x}_i$ and $\mathbf{x}_j$ and radii
$r_i$ and $r_j$. The center separation distance is

\begin{equation}
d_{ij}
=
\left\|
\mathbf{x}_j-\mathbf{x}_i
\right\|.
\end{equation}

Particle interaction occurs whenever

\begin{equation}
d_{ij}<r_i+r_j.
\end{equation}

Under these conditions the particle boundaries
occupy a common geometric region. The penetration
depth is

\begin{equation}
\delta_{ij}
=
r_i+r_j-d_{ij}.
\end{equation}

The penetration depth defines the interaction state
of the particle pair. As the particle boundaries
separate, the penetration depth continuously
decreases to zero. As interpenetration increases,
the penetration depth increases continuously.

The penetration depth therefore defines the
interaction state from which subsequent force
evaluation and state evolution are derived.


\subsubsection{Interaction Geometry}
The interaction state is described by the geometric
relationship between the source particle and the
target particle. The source orientation vector

\begin{equation}
\mathbf O_i
=
r_i\mathbf n_{ij}
\end{equation}

extends from the source center to the source
boundary along the interaction direction. The
proximity vector

\begin{equation}
\mathbf P_{ij}
=
\mathbf x_j-r_j\mathbf n_{ij}-\mathbf x_i
\end{equation}

extends from the source center to the leading edge
of the target particle.

The penetration depth is therefore

\begin{equation}
\delta_{ij}
=
\|\mathbf O_i\|
-
\|\mathbf P_{ij}\|.
\end{equation}

\subsubsection{Force Generation}

Interaction forces are generated directly from the
geometric interaction state established by the
penetration depth. Increasing penetration depth
corresponds to increasing geometric incompatibility
between the particle boundaries and therefore
produces larger interaction forces.

Each particle stores a particle-owned interaction
strength $q_i$ specified by the particle
configuration. For a particle pair, the pair
interaction strength is defined as

\begin{equation}
q_{ij}
=
\max\left(
0,
\frac{q_i+q_j}{2}
\right).
\end{equation}

The interaction force magnitude is determined from
the product of the pair interaction strength and an
interaction measure,

\begin{equation}
F_{ij}
=
q_{ij} M_{ij},
\end{equation}

where $M_{ij}$ is the interaction measure. In the
present implementation the interaction measure is
taken as the penetration depth,

\begin{equation}
M_{ij}
=
\delta_{ij},
\end{equation}

yielding

\begin{equation}
F_{ij}
=
q_{ij}\delta_{ij}.
\end{equation}

The interaction strength $q_{ij}$ controls the
magnitude of the force generated for a given
penetration depth and therefore acts as an effective
interaction stiffness between the particle pair.

The interaction strength $q_{ij}$ controls the
magnitude of the force generated for a given
penetration depth.

\begin{equation}
\mathbf{n}_{ij}
=
\frac{\mathbf{x}_j-\mathbf{x}_i}
     {\left\|
      \mathbf{x}_j-\mathbf{x}_i
      \right\|}.
\end{equation}

The force acting on particle $i$ is therefore

\begin{equation}
\mathbf{F}_{ij}
=
-F_{ij}\mathbf{n}_{ij}.
\end{equation}

\subsubsection{Acceleration and State Evolution}

The interaction force produces an acceleration that
acts along the line joining the particle centers.
The acceleration of particle $i$ is determined from

\begin{equation}
\mathbf{a}_i
=
\frac{\mathbf{F}_i}{m_i},
\end{equation}

where $m_i$ is the particle mass and
$\mathbf{F}_i$ is the accumulated source-owned force
acting on the particle.

The acceleration acts opposite the direction of
penetration and therefore reduces the component of
particle velocity directed toward the interacting
particle. As the penetration depth increases, the
resulting interaction force increases, producing a
corresponding increase in acceleration.

During the interaction, the component of particle
velocity along the interaction normal continuously
decreases until it reaches zero. Because the
interaction force remains active while penetration
exists, continued acceleration causes the normal
velocity component to change sign, resulting in
particle separation.

Consequently, rebound behavior emerges directly from
the evolution of the interaction force and does not
require explicit reflection, collision impulses, or
velocity reversal operations.

\subsubsection{Acceleration and State Evolution}

After force evaluation, the accumulated source-owned
force is converted into acceleration. For particle $i$,

\begin{equation}
\mathbf{a}_i
=
\frac{\mathbf{F}_i}{m_i},
\end{equation}

where $\mathbf{F}_i$ is the total force accumulated
by the source particle and $m_i$ is the particle mass.

Velocity is then updated from the current acceleration,

\begin{equation}
\mathbf{v}_i^{\,n+1}=\mathbf{v}_i^{\,n}+\mathbf{a}_i^{\,n}\Delta t .
\end{equation}

During penetration, the interaction force acts along
the line of interaction and opposes continued motion
into the neighboring particle. As the penetration
depth increases, the force and acceleration increase,
reducing the velocity component along the interaction
normal. When this normal velocity component reaches
zero, continued acceleration causes it to change sign.
The particle therefore transitions from compression to
rebound through the velocity update itself, without
requiring an explicit reflection or velocity-reversal
operation.

\subsubsection{Momentum Conservation}

The interaction force acts along the line joining the
particle centers and therefore produces a central
force pair. For a particle pair, the interaction
strength is symmetric,

\begin{equation}
q_{ij}=q_{ji},
\end{equation}

and the penetration depth is identical for both
particles,

\begin{equation}
\delta_{ij}=\delta_{ji}.
\end{equation}

Consequently, the pair force magnitudes satisfy

\begin{equation}
F_{ij}=F_{ji},
\end{equation}

while the interaction normals satisfy

\begin{equation}
\mathbf{n}_{ji}
=
-\mathbf{n}_{ij}.
\end{equation}

The resulting forces are therefore equal and
opposite,

\begin{equation}
\mathbf{F}_{ji}
=
-\mathbf{F}_{ij}.
\end{equation}

Because the interaction forces form a central
equal-and-opposite pair, the total linear momentum
of the particle pair is conserved.

Each particle independently evaluates the same
interaction geometry from its own source
perspective. The resulting pair forces are equal
in magnitude and opposite in direction, ensuring
momentum conservation while preserving source-owned
state updates.

Momentum conservation therefore emerges from the
symmetry of the interaction geometry and does not
require explicit pair updates or shared-state
modification.